\documentclass[11pt]{article}
\usepackage[utf8]{inputenc}
\usepackage{amsmath,amssymb}
\usepackage[margin=1in]{geometry}
\usepackage{hyperref}
\emergencystretch=3em

\title{Extending $S_\lambda$ below zero: the recursion $S_\lambda=\tfrac{\sqrt\beta}{2\lambda}bS_{\lambda+1/2}$}
\author{Claude, with Daniel Murfet}
\date{2026-09-06}

\begin{document}
\maketitle
\sloppy

\begin{abstract}
Unit 19 of the grammar-paper \S4 autoformalisation: the extension of the fluctuation function to
non-positive order (\texttt{defn:S\_negative}). For $\lambda\le0$ (away from the poles
$\{0,-\tfrac12,-1,\dots\}$) the paper defines $S_\lambda=\tfrac{\sqrt\beta}{2\lambda}\,bS_{\lambda+1/2}$,
lowering the order until the integral regime is reached. We formalise this recursion and prove the
\emph{lowering} ladder relation of \texttt{lem:extended\_ladder}, which is definitional. Zero
\texttt{sorry}/\texttt{axiom}.
\end{abstract}

\section{Setting}
The closed form $S_\lambda=2^{1-\lambda}\beta^{-\lambda}\Gamma(2\lambda)e^{\beta a^2/8}D_{-2\lambda}
(-a\sqrt{\beta/2})$ continues meromorphically in $\lambda$, but our \texttt{parCylNeg} is the
positive-order integral representation only. The paper's own extension is instead recursive, via the
lowering operator $b$ from the ladder algebra (unit~16). This is the final structural piece of \S4.

\section{The things formalised}
{\small
\begin{verbatim}
noncomputable def Sneg (β μ : ℝ) : ℕ → ℝ → ℝ
  | 0       => fun a => fluctuation β μ a
  | (k + 1) => fun a =>
      β ^ ((1:ℝ)/2) / (2 * (μ - (↑(k+1))/2)) * lowerOp β (Sneg β μ k) a

theorem Sneg_lowering (β μ : ℝ) (hβ : 0 < β) (k : ℕ) (hlam : μ - (↑(k+1))/2 ≠ 0) (a : ℝ) :
    lowerOp β (fun a => Sneg β μ k a) a
      = (2 * (μ - ↑k/2) - 1) * β ^ (-(1:ℝ)/2) * Sneg β μ (k + 1) a
\end{verbatim}
}
`Sneg β μ k` represents $S_{\mu - k/2}$, indexing the extension by the number of half-steps taken below
a base order $\mu>0$.

\section{Proof strategy}
The recursion is structural on $k$ (Lean accepts it directly; \texttt{lowerOp} uses \texttt{deriv}, which
is total, so no differentiability obligation enters the definition). The lowering relation is the
definition rearranged: $S_{\rho-1/2}=\tfrac{\sqrt\beta}{2\rho'}bS_\rho$ with $\rho'=\rho-\tfrac12$, so
$bS_\rho=\tfrac{2\rho'}{\sqrt\beta}S_{\rho-1/2}=(2\rho-1)\beta^{-1/2}S_{\rho-1/2}$. In Lean: rewrite
$S_{k+1}$ by its defining equation, group the scalar prefactor as
$\tfrac{2\rho-1}{2\rho'}\cdot\beta^{-1/2}\beta^{1/2}$, and collapse it with
$\beta^{-1/2}\beta^{1/2}=1$ and $2\rho-1=2\rho'$ (\texttt{div\_self}, needing $\rho'\ne0$ --- the
non-pole hypothesis).

\section{Roadblocks and resolutions}
Almost none: a small scalar-cancellation. Two points of care --- the numeric identity
$2(\mu-\tfrac k2)-1=2(\mu-\tfrac{k+1}2)$ is by \texttt{push\_cast; ring} (the $\uparrow(k+1)$ cast), and
the non-pole hypothesis $\mu-\tfrac{k+1}2\ne0$ is exactly what makes the prefactor invertible
(\texttt{div\_self}). (Recall \texttt{h$\lambda$} does not parse --- \texttt{$\lambda$} is the
anonymous-function token; use \texttt{hlam}.)

\section{What was Mathlib, what was new}
Only \texttt{Real.rpow\_add}/\texttt{div\_self} and \texttt{push\_cast} from Mathlib. New is the recursive
extension itself and the observation that the extended lowering relation is definitional --- no analysis
needed once the ladder operator \texttt{lowerOp} (unit~16) is in hand.

\section{Lessons}
When a special function is extended by an operator recursion, the ladder relation \emph{in the direction
of the recursion} is free (rearrange the defining equation); it is the \emph{opposite} direction (raising)
that carries the analytic content, requiring smoothness of the extension and the Weyl commutator.

\section{Follow-ups}
The remaining half of \texttt{lem:extended\_ladder} --- the raising relation
$b^\dagger S_\lambda=\beta^{1/2}S_{\lambda+1/2}$ for the extension --- and the general-$Q$ number operator
$N_\lambda S_{\lambda+Q/2}=QS_{\lambda+Q/2}$ both require the $C^\infty$-smoothness of \texttt{Sneg} (to
apply $b,b^\dagger$ and \texttt{ladder\_commutator}), proved by induction from the smoothness of the
integral $S_\mu$ and smoothness-preservation under \texttt{lowerOp}. With this, the entire ladder/log
structure of grammar \S4 is formalised.
\end{document}
